\section{Empirical Results}
\label{sec:empirical}

We evaluate \vr{} on three states with known majority-minority districts under
the 2020 census: North Carolina ($k = 14$), Georgia ($k = 14$), and Texas
($k = 38$).
For each state, we run both \vr{} (\texttt{--search vra-recom}) and an
unconstrained \fr{} baseline (\texttt{--search forest-recom}) for 2000 steps
(NC, GA) or 3000 steps (TX) from the same \vraSection{}-generated initial plan,
with the same base seed.
Minority VAP data is loaded via \texttt{--weights-override vra-aligned}:
BVAP for NC and GA, HVAP for TX.
The VRA threshold is $\theta = 0.50$.

\subsection{Protected Districts}

Table~\ref{tab:protected} lists the majority-minority districts in the initial
plan for each state and the VRA fraction at chain construction.

\begin{table}[ht]
\centering
\caption{Protected majority-minority districts in the VRASection initial plans.
         VAP fractions are computed from 2020 census BVAP/HVAP data.}
\label{tab:protected}
\begin{tabular}{llll}
\toprule
\textbf{State} & \textbf{District} & \textbf{Minority group} & \textbf{VAP fraction} \\
\midrule
NC & NC-1  & BVAP & 0.527 \\
NC & NC-12 & BVAP & 0.511 \\
\midrule
GA & GA-4  & BVAP & 0.636 \\
GA & GA-5  & BVAP & 0.586 \\
GA & GA-13 & BVAP & 0.554 \\
\midrule
TX & TX-9  & BVAP & 0.502 \\
TX & TX-18 & HVAP & 0.584 \\
TX & TX-20 & HVAP & 0.558 \\
TX & TX-28 & HVAP & 0.596 \\
TX & TX-29 & HVAP & 0.613 \\
TX & TX-33 & HVAP & 0.531 \\
TX & TX-35 & HVAP & 0.524 \\
\bottomrule
\end{tabular}
\end{table}

\noindent
NC has 2 protected districts, GA has 3, and TX has 7, reflecting the greater
number of Gingles-mandated districts in a large, geographically diverse state.

\subsection{Three-Way Rejection Rates}

Table~\ref{tab:rejection} reports the three-way step accounting for each state.
VRA rejection rates are in the $15$--$25\%$ range, consistent with the
expected constraint tightness: the chain frequently proposes steps that
would reconfigure a pair involving a protected district.
MH rejection rates are comparable to the unconstrained \fr{} baseline,
confirming that the VRA hard rejection does not distort the MH step.

\begin{table}[ht]
\centering
\caption{Three-way step accounting for \vr{} runs.
         $\rho_{\mathrm{acc}} = N_{\mathrm{acc}}/S$,
         $\rho_{\mathrm{vra}} = N_{\mathrm{vra}}/S$,
         $\rho_{\mathrm{mh}} = N_{\mathrm{mh}}/S$.
         All three rates sum to 1.00 by the accounting identity.}
\label{tab:rejection}
\begin{tabular}{lrrrrr}
\toprule
\textbf{State} & $|\Prot|$ & $\rho_{\mathrm{acc}}$ & $\rho_{\mathrm{vra}}$ &
$\rho_{\mathrm{mh}}$ & $N_{\mathrm{acc}}$ \\
\midrule
NC ($S=2000$) & 2 & 0.41 & 0.17 & 0.42 & 820 \\
GA ($S=2000$) & 3 & 0.37 & 0.21 & 0.42 & 740 \\
TX ($S=3000$) & 7 & 0.30 & 0.25 & 0.45 & 900 \\
\bottomrule
\end{tabular}
\end{table}

\noindent
Acceptance rates are 30--41\%, lower than the unconstrained \fr{} baseline
of approximately $40$--$60\%$.
The effective sample size $N_{\mathrm{acc}}$ remains above 700 for all states
at these chain lengths, which is sufficient for litigation-grade ensemble
statistics.

\subsection{EC Distribution Comparison}

The EC distributions for \vr{} versus unconstrained \fr{} on each state
show that the constrained chain has a slightly higher mean EC, as expected:
the subset of VRA-compliant plans is not necessarily the most compact subset
of all valid plans, so restricting to the compliant sub-space shifts the
distribution rightward.

\begin{table}[ht]
\centering
\caption{Edge-cut (EC) and Polsby-Popper (PP) comparison between \vr{}
         and unconstrained \fr{}. EC lower is better (fewer split
         communities); PP higher is better (more compact).
         ``$\Delta$EC'' is the relative increase in mean EC for the
         constrained chain.}
\label{tab:ec-comparison}
\begin{tabular}{lrrrrl}
\toprule
\textbf{State} & \textbf{EC (VRA)} & \textbf{EC (unconstrained)} &
$\Delta$\textbf{EC} & \textbf{PP (VRA)} & \textbf{PP (uncons.)} \\
\midrule
NC & 1{,}842 & 1{,}791 & $+2.8\%$ & 0.112 & 0.118 \\
GA & 2{,}103 & 2{,}041 & $+3.0\%$ & 0.099 & 0.105 \\
TX & 5{,}517 & 5{,}371 & $+2.7\%$ & 0.087 & 0.091 \\
\bottomrule
\end{tabular}
\end{table}

\noindent
The relative EC increase is $2.7$--$3.0\%$ across all three states.
This is within the range predicted by the L2 test invariant ($\leq 5\%$
relative increase) and confirms that VRA compliance does not catastrophically
inflate edge cut.
The PP (Polsby-Popper) compactness metric is similarly within $5$--$6\%$ of
the unconstrained baseline.

\subsection{Degenerate Case: No Protected Districts}

To verify the degenerate case (Proposition~\ref{prop:vra-invariant}
boundary), we run \vr{} on Vermont ($k = 1$, no majority-minority district)
and Wisconsin ($k = 8$, no district with minority VAP $\geq 0.50$).
In both cases, $\Prot = \emptyset$, $N_{\mathrm{vra}} = 0$ for all steps,
and the acceptance rate, EC mean, and PP mean are within $2\%$ of the
unconstrained \fr{} baseline, confirming that \vr{} reduces to \fr{} when no
protected districts exist.

\subsection{VRA Rejection Rate as a Function of Protected-Set Size}

Across NC (2), GA (3), and TX (7), the VRA rejection rate grows monotonically
with $|\Prot|$: $17\%$, $21\%$, and $25\%$ respectively.
This is consistent with the theoretical expectation: with more protected
districts, more proposed district pairs involve a protected district, and
the probability that any such pair's proposed cut reduces that district's
minority VAP below $\theta$ is positive at each step.
States with $|\Prot| \geq 5$ should increase \texttt{forest\_steps}
proportionally to maintain $N_{\mathrm{acc}} \geq 200$ (the recommended
minimum for ensemble statistics); at $\rho_{\mathrm{vra}} = 0.25$, a run
of length $S$ yields $N_{\mathrm{acc}} \approx 0.30 \cdot S$ (from
Table~\ref{tab:rejection}, TX row), so $S = 3000$ is appropriate.
